\section{Background and Related Work}\label{sec:chapter-2-background-and-related-work}

Chapter 1 stated the one-sentence thesis and the three-layer design that carries it. This chapter places that design in the literature and marks exactly where it differs from its nearest neighbours. CICADA sits at the junction of four literatures: knowledge-graph question answering (KGQA) benchmarks and semantic parsing; verifier-guided generation and self-correction; bootstrapped self-training and distillation; and LLM-over-KG agents. Its overall shape (generate candidate outputs, filter them, retrain on the survivors) is deliberately \emph{not} claimed as novel. That loop is by now a standard instrument, and several systems surveyed below share it almost exactly. The dividing line this chapter draws against the nearest neighbours runs along two axes only. The first is \textbf{filter strength}: what evidence a candidate must present before it is allowed to become supervision. This may be gold-answer agreement, a self-consistency vote, bare execution success, or, in our case, a deterministic multi-stage verifier plus a training-time-only oracle gate. The second is \textbf{blind-region handling}: what happens to the candidates that pass the filter \emph{wrongly}. They may silently poison the training pool, or they may be surfaced, labelled, and put to work. The reader should finish this chapter able to predict, for any related system, where its filter is weaker than a deterministic executable check and what that system does with its blind region.

\section{KGQA benchmarks and semantic parsing}\label{sec:kgqa-benchmarks-and-semantic-parsing}

Question answering over structured knowledge has a long semantic-parsing tradition: map a natural-language question to a formal query, execute it, return the result. The large modern KGQA benchmarks, GrailQA \citep{gu2023pangu}, LC-QuAD 2.0 \citep{dubey2019lcquad2}, and KQA Pro \citep{cao2022kqapro}, established the construction recipe this work adapts. The recipe enumerates logical forms or query templates first, verbalises them into canonical questions, then paraphrases for surface diversity, with crowd workers typically performing the verbalisation and validation stages. GrailQA also established the now-standard evaluation split taxonomy (i.i.d., compositional, zero-shot), which foregrounds whether a parser generalises beyond the query shapes it was trained on.

CICADA's benchmark (Chapter 4) adopts this template-first recipe but replaces the crowd stages with LLM stages. In particular, it replaces crowd \emph{validation} with mechanical fidelity checks. The substitution is motivated by a measured failure rather than an assumption. An LLM checker asked to validate paraphrase fidelity approved every one of 9,762 rewrites presented to it, including rewrites containing known semantic drifts (project measurement; see results table, promoted rows). An uncertain model validating an uncertain model's output detects nothing; only deterministic projection checks (entities, literals, numbers, cue words preserved) provide a signal. We turn the same scepticism on the benchmark itself. Benchmark correctness in this literature is usually asserted, whereas Chapter 4 \emph{measures} it. Every answer oracle is implemented twice from independent code paths, and their agreement rate is reported: 99.88\% (14,752/14,770) (dual-oracle audit; results table, promoted rows).

One further gap in the benchmark literature matters here. Mainstream KGQA benchmarks contain almost exclusively answerable questions. Unanswerable or ill-posed questions, and therefore abstention as a scored behaviour, are largely absent, a gap noted in work on selective prediction and answerability detection for QA \citep{rajpurkar2018squad2,kamath2020selective}. For a public-accountability domain this omission is not tolerable. Chapter 4 builds three families of abstention traps directly into the benchmark, and \S2.6 tracks abstention support as a column of the positioning table.

\section{Verifier-guided generation and self-correction}\label{sec:verifier-guided-generation-and-self-corr}

A second literature asks whether a model's own outputs can be improved after the fact. The negative result is by now well established: \emph{intrinsic} self-correction, where a model critiques and revises its own reasoning with no external signal, tends to degrade rather than improve reasoning performance \citep{huang2024intrinsic}. The positive results all share one feature: the feedback comes from an environment, and it is \emph{typed}. Self-Debugging repairs code against execution traces and error messages \citep{chen2023selfdebug}. CRITIC couples generation to external tools whose outputs ground the critique \citep{gou2024critic}. In text-to-SQL, execution-guided decoding \citep{wang2018execguided} and the decomposed prompting pipelines DIN-SQL \citep{pourreza2023dinsql} and MAC-SQL \citep{wang2024macsql} route schema and execution errors back into targeted repair. QueryAgent's ERASER performs stepwise self-correction from typed environment feedback in KGQA \citep{huang2024queryagent}. The prescription that emerges is that repair works when the error is diagnosed by something other than the model, and the diagnosis names the error type.

CICADA's \emph{reflector} follows this typed-error prescription exactly, with one positioning rule that this dissertation maintains strictly: the reflector is a \emph{consumer} of verifier signal, never a verifier itself. The four deterministic checks (Chapter 5) produce diagnostic signals: compilation failures, grounding failures, structural mismatches, and, at harvest time, oracle disagreement. The reflector, an LLM and therefore uncertain, consumes those diagnostics to propose a repaired plan, and the repaired plan re-enters the same four checks. Nothing the reflector says is ever trusted on its own authority. This design decision keeps the architecture out of the "uncertain model verifies uncertain model" regress of \S2.1. The reflector also adds one repair behaviour with no direct analogue in the text-to-SQL repair literature: provenance-gated empty-result handling. An empty result whose filter literals all trace verifiably to the question is an \emph{answer} (there genuinely are no such contracts), not a defect to be repaired away. Only an empty result with untraceable literals is a repair trigger. A development-slice measurement motivates the distinction. A naive "empty means broken, so retry" repair loop was net-negative on the \textasciitilde{}50-question dev\_smoke slice: accuracy dropped from 84\% to 66\%, and hallucinated answers rose from 1 to 6 (a development-slice measurement, restated with its gate in Chapter 5, \S5.5; motivation only, never a comparative claim).

\section{Self-training, rejection sampling, and distillation}\label{sec:self-training-rejection-sampling-and-dis}

This is the section in which the dividing line must be drawn, because the loop shape here is genuinely shared.

\subsection{Filtered self-training loops}\label{sec:filtered-self-training-loops}

STaR \citep{zelikman2022star} established the modern pattern: sample rationales from the model, keep those whose final answer matches the gold label, fine-tune on the survivors, iterate. Rejection-sampling fine-tuning generalises the pattern. RFT scales it for mathematical reasoning \citep{yuan2023rft}, and RAFT frames it as reward-ranked filtering \citep{dong2023raft}. ReST \citep{gulcehre2023rest} and ReST-EM \citep{singh2024restem} cast the grow-filter-train cycle as expectation-maximisation, and show that a model can improve on problems whose solutions it filtered itself. All of these accept a model output into the training pool on a \emph{pass criterion}. In all of them the pass criterion is the filter's entire theory of correctness: a gold final answer, a reward-model score, or a binary success signal. None of them examines closely the population of outputs that satisfy the criterion for the wrong reasons. The right answer reached by an invalid rationale in STaR is the canonical example. In all of them that population enters the training pool indistinguishably from the genuinely correct one.

\subsection{Nearest neighbours: ExeSQL and SCD}\label{sec:nearest-neighbours-exesql-and-scd}

Two recent systems are close enough to CICADA that the boundary must be drawn with care. Both deserve to be treated as serious neighbours rather than foils.

\textbf{ExeSQL} [ExeSQL; EMNLP 2025 Findings] bootstraps text-to-SQL models across SQL dialects using execution-guided rejection sampling. Candidate queries are filtered by whether they execute, the survivors form supervised data, and iterative DPO uses execution-failing queries as preference negatives. Of the published work we are aware of, this is the loop most similar in shape to Chapter 6: an executable environment supplies the filter, and the filtered pool drives both supervised and preference-based training of a local model. \textbf{SCD} [SCD; arXiv:2511.07998] is a two-stage self-consistency distillation pipeline for structured-data question answering. A teacher-distillation stage is followed by a self-distillation stage, with self-consistency voting supplying the filter. It shows, as we do, that 8B-scale students can reach strong structured-QA performance without gold logical forms.

The boundary against both is exactly the two axes announced above. First, filter strength. ExeSQL's filter is execution success; SCD's is agreement among samples. Both are correctness \emph{proxies}: a SQL query can execute cleanly and compute the wrong thing, and a model can be consistently wrong. CICADA's filter is a deterministic compile--ground--execute--verify chain with closed slot enumerations at every projectable point, plus an oracle-agreement gate at training time. The oracle gates but never authors: it discards candidates, and contributes no tokens to any training example. We do not claim that this makes the filter complete; we claim, and measure, precisely where it is incomplete. Second, blind-region handling. Because prior filters treat their pass criterion as truth, the passing-but-wrong population is invisible to them by construction. We characterise that population directly. During self-harvest on bridge-type questions, 86.7\% of generated plans passed the deterministic verifier while only 52.3\% agreed with the oracle, a 34.4-point verifier-blind stratum on this question type (self-harvest audit, n=1,351; results table, promoted rows). That stratum is exactly the failure mode execution-success filtering cannot see, and it exists at scale.

\subsection{The divide-line}\label{sec:the-divide-line}

The contribution of this dissertation relative to this literature is therefore \emph{not} the generate--filter--retrain loop. It is two things. The first is the filter: a deterministic, executable, multi-stage compile--ground--execute--verify chain with closed slot enumerations, plus a training-time oracle that gates and never authors. The second is the blind-region handling: verifier-passing-but-wrong outputs are not left to poison the pool silently. They are routed into \emph{labelled hard negatives} for preference learning, while correct abstentions become abstention supervision rather than being discarded as non-answers. On this foundation, Chapter 6 shows that a precise-but-partial-coverage filter suffices to bootstrap \emph{both} stages of a neuro-symbolic stack, the understanding stage and the planning stage, past the cloud teachers that initialised them. Chapters 4--5 show the residual blind region being closed by abstention rather than ignored. Where ExeSQL asks "does it run?" and SCD asks "does it agree with itself?", this dissertation asks which properties can be checked deterministically, what leaks past those checks, and what the leak should be used for.

\subsection{Preference optimisation and its pathologies}\label{sec:preference-optimisation-and-its-patholog}

One consequence of harvesting hard negatives is that preference optimisation becomes available. Chapter 7 reports that using near-miss hard negatives in DPO \citep{rafailov2023dpo} can be actively harmful. The pathology has recent theoretical grounding. Razin et al. [2024] show that DPO can produce \emph{likelihood displacement}, in which the probability of the preferred response falls alongside the dispreferred one. The effect is most severe when the pair is similar, which is precisely the regime of verifier-passing near-misses that differ from a correct plan by one slot or one relation. One proposed remedy is DPOP, which adds a penalty preventing the positive's likelihood from falling \citep{pal2024smaug}. Another is IPO, which replaces DPO's unbounded objective with a bounded one that resists overfitting to the preference margin \citep{azar2024ipo}. This dissertation engages with this literature empirically rather than theoretically. Chapter 7 uses it as the interpretive frame for a mechanism-explained negative result contrasting subtractive preference suppression with additive distillation, and evaluates an IPO variant among the mitigations.

\section{LLM-over-KG agents}\label{sec:llm-over-kg-agents}

A parallel line couples LLMs to knowledge graphs at inference time. Pangu \citep{gu2023pangu} keeps generation on a leash: the LLM \emph{discriminates} among candidate logical-form extensions enumerated from the graph, rather than generating freely. This guarantees groundedness by construction, at the cost of a search procedure in the loop. ChatKBQA \citep{luo2024chatkbqa} has a fine-tuned LLM generate a logical form and then grounds its entities and relations by retrieval, decoupling generation from grounding. RoG \citep{luo2024rog} plans relation paths that are grounded against the graph and then reasons over the retrieved paths. StructGPT \citep{jiang2023structgpt} gives a frozen LLM iterative read access to structured data through specialised interfaces.

CICADA shares this family's premise, that the graph and not the model is the ground, but differs in where authority sits and in what failure produces. First, the KG executor is the \emph{sole answer authority}. The LLM never composes the answer text from retrieved material; it produces a plan, and only the deterministic execution of that plan produces an answer. Second, the plan is a \emph{checkable artifact}: a typed graph plan compiled by a deterministic twelve-transform compiler, so conformance to the executable contract is a mechanically decidable property rather than a hope about generation. Third, grounding failures are \emph{diagnosable events} rather than dead ends. A failed grounding emits a typed diagnostic that feeds the gated repair of \S2.2 or, when repair is exhausted, a principled abstention. Where Pangu prevents ungrounded output by restricting the model to discrimination, CICADA permits generation and then subjects it to checks whose failure modes are themselves informative. This design difference is what makes the harvest of \S2.3 possible in the first place.

\section{LLM-generated evaluation data}\label{sec:llm-generated-evaluation-data}

Using LLMs to generate benchmarks invites a known pathology: naive paraphrasing collapses toward a house style, which overstates surface diversity while narrowing it [see e.g. Yu et al. 2023, AttrPrompt]. Chapter 4's generation pipeline conditions paraphrasing on six explicit style axes with mechanical acceptance checks. In keeping with this chapter's theme of measuring what is usually asserted, it reports distributional evidence rather than a diversity claim: the median pairwise trigram-Jaccard similarity among accepted surfaces is 0.429 (distributional measurement; results table, promoted rows). Two disciplinary notes bound the claim. First, we do not assert that these LLM-generated surfaces are more "natural" than crowd-written ones. The opposite concern, that they are less natural, is a documented limitation taken up in Chapter 8; the claim is measured correctness and measured diversity, nothing more. Second, paraphrase must preserve \emph{abstention cues}. A rewrite that smooths away the ambiguity that made a question unanswerable silently converts an abstention item into an answerable one. Cue preservation is therefore enforced as a benchmark-integrity device (Chapter 4).

\section{Positioning summary}\label{sec:positioning-summary}

Table 2.1 operationalises the divide-line. Each neighbour is scored on the two axes of this chapter, filter strength and blind-region handling, together with abstention support. No column concerns loop shape, because loop shape is not where the difference lies. No accuracy numbers appear, because numbers measured on different benchmarks are not comparable, and none of this dissertation's numbers were measured on theirs.

\textbf{Table 2.1 --- Positioning against neighbour systems.}

\begin{table}[t]
\centering
\small
\caption{Positioning against the nearest neighbours: filter strength, blind-region handling, stages distilled, and abstention support.}
\label{tab:02_background_1}
\begin{tabular}{lllll}
\toprule
System & Filter type & Filter strength & Blind-region handling & Abstention support \\
\midrule
STaR \citep{zelikman2022star} & gold final-answer match & gold-dependent; rationale unchecked & right-answer/wrong-rationale enters pool silently & none \\
RFT / RAFT \citep{yuan2023rft} & answer match / reward-model score & proxy reward; unexecutable & reward-hacking outputs indistinguishable & none \\
ReST / ReST-EM \citep{gulcehre2023rest,singh2024restem} & scalar reward threshold & proxy reward & above-threshold-but-wrong enters pool & none \\
ExeSQL [EMNLP 2025 Findings] & execution success (+ iterative DPO on execution failures) & executable but success-only & executes-but-wrong passes; negatives are execution \emph{failures} (easy negatives) & none \\
SCD [arXiv:2511.07998] & self-consistency vote (after teacher stage) & agreement, not correctness & consistently-wrong passes silently & none \\
Pangu \citep{gu2023pangu} & grounded enumeration (inference-time) & strong but generation-restricting & n/a (no training pool) & none \\
ChatKBQA \citep{luo2024chatkbqa} / RoG \citep{luo2024rog} & retrieval grounding of generated forms/paths & grounding only; trained on gold & grounding failures not supervision & none \\
\textbf{CICADA} & deterministic compile--ground--execute--verify, closed slot enums; training-time oracle gate (filters, never authors) & multi-stage, executable, deterministic; coverage partial and measured & verifier-passing-but-wrong $\rightarrow$ labelled hard negatives; correct abstentions $\rightarrow$ abstention supervision & first-class output, three trap families \\
\bottomrule
\end{tabular}
\end{table}


Two boundary clarifications keep the table honest. The oracle-agreement gate exists \textbf{only during training-data harvest}. At evaluation time no oracle signal enters the pipeline, since that would constitute answer leakage, so the inference-time system is filtered exactly as strongly as its deterministic checks, no more. And the reflector does not appear in the filter column anywhere: it consumes the filter's diagnostics, and its outputs re-enter the same checks, so it adds repair capacity, not filter strength.

With the neighbours placed, the remaining chapters proceed constructively. Chapter 3 describes the data, and Chapter 4 presents the benchmark whose integrity apparatus \S2.1 motivated. Chapter 5 presents the three-layer system whose checks define the verifiable region. Chapter 6 presents the bootstrap experiments that draw supervision from it, and Chapter 7 analyses what leaks past it.
